\section*{PHẦN KẾT LUẬN}
\addcontentsline{toc}{section}{PHẦN KẾT LUẬN}

Vận dụng thế giới quan và phương pháp luận của chủ nghĩa duy vật biện chứng và chủ nghĩa duy vật lịch sử để nghiên cứu đề tài ``Tìm hiểu về tầng lớp nho sĩ ở Việt Nam thế kỷ XVIII'', tiểu luận rút ra ba kết luận chủ yếu sau đây:

Thứ nhất, sự biến đổi của tầng lớp nho sĩ thế kỷ XVIII là minh chứng điển hình cho quy luật tồn tại xã hội quyết định ý thức xã hội và vị trí trung gian đặc thù của trí thức trong xã hội có giai cấp. Cuộc khủng hoảng toàn diện của thể chế Lê -- Trịnh, chúa Nguyễn, sự thương mại hóa khoa cử với hiện tượng ``Sinh đồ ba quan'' năm 1750, sự khởi sắc của kinh tế hàng hóa cùng làn sóng khởi nghĩa nông dân mà đỉnh cao là phong trào Tây Sơn đã làm sụp đổ thế độc tôn của Nho giáo chính thống. Chính sự biến động từ gốc rễ của cơ sở kinh tế -- xã hội đã buộc người nho sĩ phải xét lại các tín điều cũ để lựa chọn con đường hành xử mới.

Thứ hai, tầng lớp nho sĩ thế kỷ XVIII đã diễn ra sự phân hóa quyết liệt về lập trường chính trị, song hành cùng bước chuyển mình sâu rộng về học thuật và tư tưởng nhân văn. Trên bình diện chính trị, khác với bộ phận bảo thủ trói mình trong chữ ``trung quân'' giáo điều và bộ phận bất mãn lui về nông thôn mở trường tư thục giữ gìn văn hiến, những đại biểu ưu tú nhất như Ngô Thì Nhậm, Nguyễn Thiếp, Phan Huy Ích đã vượt qua rào cản đẳng cấp để dấn thân phò tá vua Quang Trung bảo vệ độc lập dân tộc; còn ở phương Nam, danh sư Võ Trường Toản với phương châm ``Tri ngôn dưỡng khí'' đã đào tạo nên lớp hiền tài làm rạng danh sĩ khí Gia Định. Trên bình diện học thuật, kẻ sĩ thế kỷ XVIII đã vượt thoát lối học từ chương Tống Nho để kiến tạo trào lưu thực học bách khoa mà đỉnh cao là tư tưởng ``Lý ở trong Khí'' và phương pháp khảo chứng thực địa của Lê Quý Đôn, kết hợp với xu hướng dung hợp Tam giáo và sự bùng nổ của chủ nghĩa nhân đạo trong văn học chữ Nôm.

Thứ ba, nhìn nhận khách quan theo nguyên tắc lịch sử -- cụ thể, dù chưa thể vượt qua giới hạn của phương thức sản xuất phong kiến tiểu nông và khuôn khổ lý tưởng ``minh quân -- nhân chính'', tầng lớp nho sĩ thế kỷ XVIII đã để lại di sản quý giá về tinh thần thực học, ý thức chủ quyền quốc gia và khát vọng giải phóng con người. Những kinh nghiệm lịch sử ấy mang ý nghĩa tham chiếu trực tiếp đối với hiện tại: kiên quyết bài trừ bệnh hư danh và bằng cấp hình thức, gắn kết chặt chẽ nghiên cứu lý luận với thực tiễn phát triển đất nước, bồi đắp bản lĩnh chính trị của người trí thức và thực thi hiệu quả chính sách trọng dụng hiền tài theo tinh thần Nghị quyết số 45-NQ/TW của Ban Chấp hành Trung ương Đảng Khóa XIII trong giai đoạn phát triển mới.
